\section*{Capítulo \ref{ch:b3:genetic-engineering} --- Engenharia genética e biotecnologia}

\begin{solution}{exo:b3:genetic-engineering:1}
Uma origem de replicação, para que o hospedeiro copie o \omterm{def:b3:genetic-engineering:tools}{plasmídeo} e o
passe às células-filhas; um \omterm{def:b3:genetic-engineering:tools}{marcador de seleção}, em geral de
resistência a antibiótico, para que só cresçam as células que carregam
o \omterm{def:b3:genetic-engineering:tools}{plasmídeo}; e um sítio de restrição único (ou um agrupamento de
sítios) em que o inserto é ligado sem cortar o \omterm{def:b3:genetic-engineering:tools}{plasmídeo} em outro
lugar.
\end{solution}

\begin{solution}{exo:b3:genetic-engineering:2}
Uma dada sequência de seis bases ocorre uma vez a cada $4^{6} = 4096$
bp em DNA aleatório: cerca de $4.6\times 10^{6}/4096 \approx 1100$
fragmentos de \qty{4.1}{kb} em média. Os genomas reais não são
aleatórios --- a composição de bases, o uso de códons e a evitação de
algumas sequências deslocam a contagem (o número real de sítios EcoRI
em \emph{E.~coli} é de várias centenas).
\end{solution}

\begin{solution}{exo:b3:genetic-engineering:3}
A desnaturação a \qty{95}{\celsius} separa as fitas; o anelamento a
\qtyrange{50}{65}{\celsius} deixa os \omterm{def:b3:genetic-engineering:pcr}{iniciadores} parearem com seus
sítios; a extensão a \qty{72}{\celsius} deixa a polimerase copiar a
partir de cada \omterm{def:b3:genetic-engineering:pcr}{iniciador}. Uma polimerase comum seria destruída a
\qty{95}{\celsius} já no primeiro ciclo e teria de ser acrescentada de
novo trinta vezes; a Taq sobrevive a todos os ciclos.
\end{solution}

\begin{solution}{exo:b3:genetic-engineering:4}
Um \omterm{def:b3:genetic-engineering:crispr}{RNA-guia} cujas vinte bases correspondam ao alvo, e um \omterm{def:b3:genetic-engineering:crispr}{PAM} (NGG)
imediatamente a 3$'$ do alvo na fita não alvo. Depois da \omterm{def:b3:dna-repair:lesions}{quebra de fita
dupla}: a junção de extremidades, que deixa uma pequena inserção ou
deleção e em geral destrói a fase de leitura do gene; ou a \omterm{def:b3:dna-repair:dsb}{recombinação
homóloga} com um molde que carrega a sequência desejada, que a escreve
ali.
\end{solution}

\begin{solution}{exo:b3:genetic-engineering:5}
$50\times 1.9^{35} = 50\times 5.7\times 10^{9} \approx 3\times 10^{11}$
cópias. $\Delta C_{T} = 7.3$; razão $1.95^{7.3} = e^{7.3\ln 1.95}
\approx 130$: a primeira amostra tinha cerca de $130$ vezes mais molde.
\end{solution}

\begin{solution}{exo:b3:genetic-engineering:6}
As bactérias não fazem splicing: uma cópia genômica com íntrons seria
transcrita num mensageiro inutilizável, e por isso se usa o cDNA sem
íntrons, copiado do mRNA maduro. Outros obstáculos: o promotor humano
não é reconhecido pela polimerase bacteriana (forneça um promotor
bacteriano e um sítio de ligação ao ribossomo); o uso de códons humano
é diferente (sintetize o gene com os códons preferidos do hospedeiro);
a proteína pode não se enovelar ou pode se agregar (baixe a
temperatura, funda-a a um parceiro solúvel ou reenovele-a a partir dos
corpos de inclusão); a glicosilação está ausente (use levedura ou
células de mamífero se ela importar); as pontes dissulfeto precisam do
periplasma oxidante.
\end{solution}

\begin{solution}{exo:b3:genetic-engineering:7}
As células-tronco direcionadas são injetadas num blastocisto
hospedeiro e se misturam com as células dele, de modo que o camundongo
é construído a partir de dois genótipos --- uma quimera --- e só
transmite o alelo se algumas de suas células germinativas derivarem das
células direcionadas. Cruzar a quimera com um camundongo selvagem dá
heterozigotos (vindos de metade de seus gametas, se metade da linhagem
germinativa for direcionada); intercruzar heterozigotos dá um quarto de
\omterm{def:b3:genetic-engineering:transgenic}{nocautes} homozigotos.
\end{solution}

\begin{solution}{exo:b3:genetic-engineering:8}
Exatas: $6.4\times 10^{9}\times 4^{-20}\times (1/16) \approx 4\times
10^{-4}$ sítio (o GG do \omterm{def:b3:genetic-engineering:crispr}{PAM} custa um fator $16$) --- nenhum, na
prática. A até duas discordâncias há $1 + 60 + 1710 = 1771$ sequências:
$1771\times 4\times 10^{-4} \approx 0.7$ sítio de acaso. O genoma real
não é aleatório, e um guia é conferido diretamente contra ele.
\end{solution}

\begin{solution}{exo:b3:genetic-engineering:9}
$p_{1} = 0.05 + 0.8\times 0.05\times 0.95 = 0.088$; $p_{2} = 0.152$;
$p_{3} = 0.255$; depois $0.41$, $0.60$: cerca de cinco gerações para
passar de $0.5$, contra a estimativa
$\ln(1/(2\times 0.05))/\ln 1.8 \approx 4$.
\end{solution}

\begin{solution}{exo:b3:genetic-engineering:10}
A junção de extremidades age em cada célula, inclusive nas
células-tronco quiescentes, em que a \omterm{def:b3:dna-repair:dsb}{recombinação homóloga}, um processo
de S/G2, é ineficiente; ela não exige a entrega de um molde e seu
produto --- qualquer ruptura do intensificador --- resolve o problema,
ao passo que a correção exige a sequência exata e rende uma minoria de
células. Além disso, a mesma edição serve para qualquer mutação da
$\beta$-globina, falciforme ou talassêmica. Desvantagem: ela remove um
elemento regulatório (com quaisquer outros papéis que ele tenha), deixa
o alelo falciforme no lugar e a mistura de indels fica sem controle.
\end{solution}

\begin{solution}{exo:b3:genetic-engineering:11}
Com adequação $1 - s$ nos homozigotos: $p' = \bigl[p^{2}(1-s) + p(1-p)(1+c)
\bigr]/(1 - sp^{2})$. A partir da raridade, os homozigotos
são desprezíveis e o impulsor cresce como $(1+c)p$ seja qual for $s$;
se ele vai à fixação decide-se perto de $p = 1$, onde um alelo
selvagem raro de frequência $q$ retorna como
$q' \approx q(1-c)/(1-s)$: há fixação se $c > s$, e um equilíbrio
intermediário em caso contrário. Resistência: cada conversão que falha
(junção de extremidades em vez de recombinação) deixa um alelo que o
guia já não reconhece; se esses alelos forem aptos --- um indel em fase
numa região não essencial --- eles não têm custo algum enquanto o
impulsor tem, de modo que a seleção os favorece e eles substituem o
impulsor. Os remédios são mirar sequências essenciais em que os indels
sejam letais e usar vários guias ao mesmo tempo.
\end{solution}

\begin{solution}{exo:b3:genetic-engineering:12}
A \omterm{def:b3:genetic-engineering:crispr}{Cas9} injetada num zigoto pode agir depois das primeiras divisões, de
modo que células diferentes recebem reparos diferentes --- o mosaicismo;
o reparo de cada quebra é escolhido pela célula, com a junção de
extremidades dando indels imprevisíveis e deleções grandes ou perda de
heterozigosidade no corte; o desfecho por \omterm{def:b3:dna-repair:dsb}{recombinação homóloga} é
minoritário; quebras \omterm{prop:b3:genetic-engineering:editors}{fora do alvo} ocorrem em sítios que não se pode
prever todos; e o embrião só pode ser conferido sequenciando algumas
células biopsiadas, que não podem falar pelas demais. Evidência
necessária: métodos que editem cada célula de modo idêntico (antes da
primeira fase S) com eficiências próximas de \qty{100}{\%},
sequenciamento do genoma inteiro de cada célula de embriões animais
editados mostrando nenhuma mudança não intencional, e acompanhamento de
longo prazo de animais editados ao longo de gerações --- nada do que
existe hoje.
\end{solution}

\begin{solution}{pb:b3:genetic-engineering:1}
\textbf{1.} Uma vez a cada $4096$ bp. $P(\text{sem sítio}) = (1 - 1/4096)^{1500}
= e^{-0.37} = 0.69$.
\textbf{2.} Escolha uma enzima sem sítio no gene; ou amplifique o gene
por PCR com \omterm{def:b3:genetic-engineering:pcr}{iniciadores} que carreguem novos sítios de restrição em suas
extremidades 5$'$ (ou use clonagem de extremidades cegas ou
independente de ligação).
\textbf{3.} \qty{0.1}{\micro g}\,$\times 2\times 10^{7} = 2\times
10^{6}$ colônias; \qty{10}{\%}, ou $2\times 10^{5}$, carregam o
inserto.
\textbf{4.} Brancas: \qty{10}{\%}. Metade dos insertos está na
orientação certa, de modo que $P(\text{nenhum em } n \text{ brancas}) = 0.5^{n} \le 0.01$ dá $n \ge 6.6$: pesque sete.
\textbf{5.} $32$ duplicações, $2^{32} = 4.3\times 10^{9}$ células,
$2.1\times
10^{11}$ \omterm{def:b3:genetic-engineering:tools}{plasmídeos}; cada um com $5500\times 650 = 3.6\times 10^{6}$ Da $=
5.9\times 10^{-18}$ g: \qty{1.3}{\micro g} de
\omterm{def:b3:genetic-engineering:tools}{plasmídeo}.
\textbf{6.} A clonagem só precisa de replicação, fornecida pela origem
do \omterm{def:b3:genetic-engineering:tools}{plasmídeo}; o promotor humano não é lido pela RNA-polimerase
bacteriana, de modo que, para a expressão, um promotor bacteriano e um
sítio de ligação ao ribossomo precisam ser postos à frente da sequência
codificadora (sem íntrons).
\textbf{7.} $N_{0} = N_{T}/2^{C_{T}}$: $10^{12}/2^{28} \approx 3700$
cópias; em $C_{T} = 35$, $10^{12}/2^{35} \approx 29$ cópias.
\textbf{8.} $2^{n} = 10^{12}$: $n = 39.9$, quarenta ciclos.
\textbf{9.} $10^{12}/2^{38} \approx 3.6$ cópias. Além de uns $40$
ciclos, uma única molécula contaminante, ou artefatos de \omterm{def:b3:genetic-engineering:pcr}{iniciador},
alcançam o limiar, e uma amostra de menos de uma cópia não significa
nada.
\textbf{10.} Uma molécula cruza no ciclo $40$: um positivo. Impeça isso
com salas e equipamentos separados para a montagem e para a manipulação
dos produtos, fluxo de trabalho de mão única, controles negativos em
toda corrida e destruição enzimática dos amplicons arrastados.
\textbf{11.} $3700/0.4 \approx 9000$ cópias de RNA.
\textbf{12.} Relatado $2^{3.3} \approx 10$ vezes; verdadeiro
$1.9^{3.3} =
e^{3.3\ln 1.9} \approx 8.3$ vezes.
\textbf{13.} $40\times \qty{35}{\micro g} = \qty{1.4}{mg}$ por dia,
\qty{0.51}{g} por ano; $\times 10^{7}$: \qty{5100}{kg} por ano.
\textbf{14.} $5.1\times 10^{6}/5800 \approx 880$ mol; $5.3\times 10^{26}$
moléculas.
\textbf{15.} $5.1\times 10^{6}$ g $/ (20$ g/L$) = 2.6\times 10^{5}$ L
$= \qty{260}{m^3}$ --- um décimo da piscina.
\textbf{16.} $5100\times 8000 = 4.1\times 10^{7}$ kg de pâncreas, a
\qty{0.1}{kg} cada: $4\times 10^{8}$ porcos por ano.
\textbf{17.} O hormônio recombinante é idêntico ao humano, e por isso
não suscita anticorpos nem causa reações alérgicas; não carrega
patógenos animais (vírus, príons); o abastecimento não depende do
abate; e a sequência pode ser aprimorada.
\textbf{18.} Os resíduos que tocam o receptor ficam inalterados, de modo
que a ligação e a sinalização são as da insulina; os resíduos alterados
ficam na superfície em que as moléculas de insulina se associam em
dímeros e hexâmeros, e mudá-los muda a rapidez com que o depósito
injetado se desfaz em monômeros absorvíveis.
\textbf{19.} $6.4\times 10^{9}\times 4^{-20}/16 \approx 3.6\times
10^{-4}$: nenhuma correspondência exata de acaso.
\textbf{20.} $1 + 60 + 1710 + 27\times 1140 = \num{32551}$ sequências;
$\num{32551}\times 3.6\times 10^{-4} \approx 12$ sítios \omterm{prop:b3:genetic-engineering:editors}{fora do alvo} de
acaso.
\textbf{21.} $10^{7}\times 10^{-4} = 1000$ células. Uma célula-tronco
vive e se divide por toda a vida do paciente; uma que tenha perdido um
supressor de tumor, ou adquirido uma translocação, pode fundar um clone
que se torna uma leucemia.
\textbf{22.} $p_{1} = 0.0019$, $p_{2} = 0.0036$, $p_{3} = 0.0069$; um
crescimento geométrico de $1.9$ por geração dá $\ln 500/\ln 1.9
\approx 10$; iterando a recorrência, a frequência passa de $0.5$ na
geração $11$ ($0.45 \to 0.68$).
\textbf{23.} Cerca de um ano. Um custo de \qty{20}{\%} nos homozigotos é
menor que $c = 0.9$, de modo que o impulsor ainda vai à fixação, mais
devagar, enquanto a adequação média da população cai --- que é
justamente o propósito de um impulsor de supressão; um custo acima de
$c$ o deteria numa frequência intermediária.
\textbf{24.} $10^{6}\times 0.1 = 10^{5}$ alelos de resistência em uma
geração. Sendo incortáveis e, se aptos, sem custo enquanto o impulsor
tem custo, eles são favorecidos e se propagam, e o impulsor é
eliminado --- por isso os impulsores miram sítios em que os produtos de
junção de extremidades sejam letais, com vários guias ao mesmo tempo.
\textbf{25.} Cerca de $3700$ cópias na primeira amostra; uns
\qty{260}{m^3} de fermentador por ano para a insulina do mundo; cerca
de $11$ gerações, um ano de mosquito, para o impulsor.
\end{solution}
